\section{Numerical experiments}

\subsection{Results for Darcy's problem}

\subsubsection{Illustrative examples}

As a continuation of the numerical experiment shown in the previous section, we aim to 
qualitatively illustrate the effect of the data set size on the data-driven solution. To that end, we first fix the number of points $N_g$ to the value $15$ from the 1D uniaxial test and rotate the corresponding
data by $N_{\theta} = 4,~6,~8,~10$ and $20$ different angles equally distributed
over $[0,2\pi)$. In this case we consider data not exhibiting deviations with
respect to the theoretical curve (the red line in figure \ref{fig:datadistproc}). 
The effect of noise in the data will be evaluated later on.
The pressure field $\uD_h$ (contours and surface) and the flux $\sD_h$ 
solutions are shown in figure \ref{fig:sol_varNtheta}. It is also interesting
to see the distribution of data over the domain $\Omega$ shown in the last
row. It is worth mentioning that, for the proximity measure chosen, the corresponding
data-driven problem becomes a linear one, whereas as, the finite element problem
solve to find the reference solution is non-linear in such case.  
%We also show the dual fields $\lD_h$ and $\mD_h$ in \ref{fig:dualsol_varNtheta}.
%On the other hand, in figure \ref{fig:sol_varNg} we show the solution when 
%$N_{\theta} = 40$ is fixed and the number of points in the 1D uniaxial test is 
%varied, namely, $N_g = 5,~10,~15,~20,~25,~30,~35$. 
The first striking observation is that the data-driven solution becomes closer to the reference one as the total number of points in the data set increases. Later in this section, we present quantitative numerical evidence of this convergence for a problem with a manufactured solution. Additionally, note that if we assume a different model to compute the reference solution, the data-driven approach will yield a different result. For example, if we assume a much simpler linear relationship between flux and gradient, both the reference and, consequently, the data-driven solution will change. This is illustrated in figure \ref{fig:sol_darcysimple}
for the same set of values $N_{\theta}$ of the previous experiment.
Another aspect that influenced the data driven solution is the fact that 
isotropy was assumed to compute the reference solution used
to initially distribute the data. If we consider otherwise an anisotropic behavior
we end up with different results. To that end we produce synthetic results
using Forcheimer's model with different permeability values assigned to two mutually 
orthogonal directions. The data points relating flux and gradient for each 
direction are shown in figure \ref{fig:anydatapoints}

\begin{figure}[ht!!]
	\begin{center}
		\caption{Data-Driven solutions produced for different values of
		$N_{\theta}$ corresponding to the Forcheimer's model
	to compute the reference solution.} \label{fig:effectNtheta}
		\fontsize{5pt}{6pt} \selectfont
		\def\svgwidth{12.5cm}
		\input{Figures/pressures_and_fluxes_forcheimer_all.pdf_tex}
	\end{center}
\end{figure}

\subsubsection{Convergence to a manufactured solution}

We assess the convergence of the data-driven formulation and all its 
variants, namely, formulation DNLSP1-A, Aeqor, B, C and D to a manufactured solution proposed as the reference case. To that end we choose $ \uD \in \mathcal{C}^{\infty} $ to be given by:
\begin{equation}
\uD:\Omega \to \mathbb{R},~p(x,y) = \cos(\pi x)\cos(\pi y), \label{eq:manup}
\end{equation}
where the domain $\Omega \subset \mathbb{R}^{2} \mid \Omega = [0,1]\times[0,1]$. Assuming the permeability to be isotropic with $\frac{k}{\mu} = 1$ we derive from (\ref{eq:manup})  the complete solution and corresponding source term to be
\begin{equation*}
	\begin{aligned}
		&\sD: \Omega \to \mathbb{R}^{2}, & &\sD = [-\pi\sin(\pi x)\cos(\pi y), -\pi \cos(\pi x)\sin(\pi y)]^{\top} \\
		&\eD: \Omega \to \mathbb{R}^{2}, & &\eD = [-\pi\sin(\pi x)\cos(\pi y), -\pi \cos(\pi x)\sin(\pi y)]^{\top} \\
		&q: \Omega \to \mathbb{R}, & &q = -2\pi^{2}\cos(\pi x) \cos(\pi y),
	\end{aligned}
\end{equation*}

First, we assess the behavior of original formulation DNLSP1-A. So, we introduce a finite element structured mesh made of $20.000$ triangular elements and produce synthetic data having different values of $N_{m} = N_{g} \times N_{\theta}$, i.e. starting with $2 \times 2$ and increasing up to $16 \times 16$. The numerical experiments, then, shown a gradual convergence of the data-driven solutions for all the unknown field -- $\uD_h$, $\sD_h$, $\eD_h$, $\lD_h$, and $\mD_h$ -- towards the manufactured solution as $ N_{m} $ increases. It is worth to mentioning that the synthetic data are produced in such way that the minimum and maximum values of the pressure gradient and the associated flux given by the manufactured solution are captured. This approach prevents mismatches between the smallest and largest physical states, $ (\eD,\sD) $, present in both reference solution and the synthetic data set, i.e. $ \mathfrak{D}_m = \left\lbrace\left(\mathfrak{g}_n,\mathfrak{s}_n\right)\right\rbrace_{n \in \mathcal{N}_m} $.


Figures \ref{fig:DNLSP1_A_plt_over_line_u}(a) to \ref{fig:DNLSP1_A_plt_over_line_u}(d), for example, present the data-driven solutions for the scalar pressure field. As can be observed, the greater the number of sampled pairs  ($N_{g} \times N_{\theta})$ the closer the data-driven solution is to the manufactured pressure field -- c.f. figure \ref{fig:DNLSP1_conv_u_L2}(f) --. Figure \ref{fig:DNLSP1_A_plt_over_line_u}(e), in addition,  provides the pressure profiles along the diagonal coordinate for each considered sampled pair relation. This further corroborates the observed convergence behavior of the data-driven solution. 

\begin{figure}[h!]
	\centering
	\caption{Figures (a) to (d) on the left present the data-driven solutions for the scalar pressure field based on DNLSP1-A formulation, using a fixed finite element mesh with 100 elements per coordinate direction and with different pairs of sampled points. Figure (e) on the right shows the data-driven solution for the pressure profiles along the diagonal coordinate corresponding to each solution depicted in figures (a) to (d).}
	\fontsize{6pt}{10pt}\selectfont 
	\def\svgwidth{\columnwidth}
	\input{Figures/DNLSP1_A_plt_over_line_u.pdf_tex}
	\label{fig:DNLSP1_A_plt_over_line_u}
\end{figure}

Figures \ref{fig:DNLSP1_A_plt_over_line_s}(a) to \ref{fig:DNLSP1_A_plt_over_line_s}(d), analogously to the previous figure, describe the data-driven solution for the flux vector field. Once again, numerical evidences shown that, as the number of sampled pairs increases, the data-driven solution progressively approaches the manufactured flux vector field -- c.f. figure \ref{fig:DNLSP1_conv_s_L2}(f) --. It should be noted that, compared to the scalar pressure field, the convergence in this case appears to be slower, since with the same choice $ N_{g}\times N_{\theta} = 16 \times 16$, the solution represented in figure \ref{fig:DNLSP1_A_plt_over_line_u}(d) is closer to the reference solution in figure \ref{fig:DNLSP1_conv_u_L2}(f) than the solution in figure \ref{fig:DNLSP1_A_plt_over_line_s}(d) is to the reference in figure \ref{fig:DNLSP1_conv_s_L2}(f). This difference in convergence rates will be justified through an convergence study of both fields in the $ L^{2} $-norm, which will be addressed in the next section.

\begin{figure}[h!]
	\centering
	\caption{Figures (a) to (d) on the left present the data-driven solutions for the vector flux field based on DNLSP1-A formulation, using a fixed finite element mesh with 100 elements per coordinate direction and with different pairs of sampled points. Figure (e) on the right shows the data-driven solution for the flux magnitude profiles along the diagonal coordinate corresponding to each solution depicted in figures (a) to (d).}
	\fontsize{6pt}{10pt}\selectfont 
	\def\svgwidth{\columnwidth}
	\input{Figures/DNLSP1_A_plt_over_line_s.pdf_tex}
	\label{fig:DNLSP1_A_plt_over_line_s}
\end{figure}

  
Figure \ref{fig:DNLSP1_A_plt_over_line_lm} finally presents the data-driven solutions for the dual fields, $ \lD_h $ and $ \mD_h $. For both cases, a behavior similar to that obtained so far is observed, that is, as the number of sampled pairs increases, the data-driven solution gradually approaches the exact solution, which in this case corresponds to $ \lD = 0 $ and $ \mD = [0,0]^{\top} $, respectively. 


\begin{figure}[h!]
	\centering
	\caption{Data-driven solutions for the dual fields based on DNLSP1-A formulation, using a fixed finite element mesh 100 elements per coordinate direction and different pairs of sampled points. Figures (a) to (d) on the left correspond to the mass balance Lagrange multiplier, $ \lD_h$, whereas figures (e) to (h) correspond to the compatibility Lagrange multiplier, $ \mD_h $.}
	\fontsize{6pt}{10pt}\selectfont 
	\def\svgwidth{\columnwidth}
	\input{Figures/DNLSP1_A_plt_over_line_lm.pdf_tex}
	\label{fig:DNLSP1_A_plt_over_line_lm}
\end{figure}


Furthermore, similar studies have been performed considering the variants of the standard data-driven formulation, namely, DNLSP1-Aeqo, B, C, and D. However, since all of them showed similar results for all unknown fields, we chose to omit them here and discuss only their convergence aspects in the next section.

One could now inquire about the accuracy of the data-driven solution for the proposed problem when the number of sampled points is fixed and the finite element mesh is progressively refined. To address this question, we conducted a numerical study for the standard formulation, DNLSP1-A, and all its variants. In this study, we fixed the number of sampled points at \( N_{\theta} = 6 \times 6 \) and gradually refined the finite element mesh, starting from a coarse grid (e.g., \( N_{x} \times N_{y} = 10 \times 10 \) elements per coordinate direction) to a fine grid (e.g., \( N_{x} \times N_{y} = 100 \times 100 \) elements per coordinate direction). 

\begin{figure}[h!]
	\centering
	\caption{Data-driven solutions for the scalar pressure field for the natural formulation and all of its variants with a fixed number of sampled points $ N_{g} \times N_{\theta} = 6 \times 6 $ and three different number of finite element meshes. }
	\def\svgwidth{\columnwidth}
	\input{Figures/sol_DNLSP1_plt_over_line_u_fem_mesh.pdf_tex}
	\label{fig:sol_DNLSP1_plt_over_line_u_fem_mesh}
\end{figure}

Figure \ref{fig:sol_DNLSP1_plt_over_line_u_fem_mesh} presents the results of that study for the scalar pressure field, \( \uD_h \). In the lower left corner of each subfigure, we show the two-dimensional pressure scalar field, ranging from the coarse mesh to the fine mesh in a descending order. Additionally, to better visibility, we plotted the solution profiles along the diagonal coordinate for each mesh and compared them to the reference or exact solution.  We first observe that the data-driven solutions based on formulations A, Aeqo, and D produce closely matching results, which is a natural consequence of the selected finite element space for that unknown. The same conclusion applies to the data-driven solutions based on formulations B and C. The staircase behavior of the numerical solution is particularly evident for the latter formulations, especially on coarse meshes. Although different characteristics are observed in each data-driven formulation, they all gradually converge to the exact solution as the finite element mesh becomes finer. However, the convergence is asymptotic up to a certain threshold, with the error remaining above a certain nonzero limit, since it depends simultaneously on the number of sampled points and the refinement of the finite element mesh.


%This can be verified in table \ref{tab:DNLSP1_fem_mesh_errors}. For formulation A, for instance, we observe that doubling the number of finite elements per coordinate direction causes the value of \( \| \uD - \uD_h \|_{L^{2}} \) to change only in the fourth decimal place, with similar behavior observed for the other formulations.


\begin{figure}[h!]
	\centering
	\caption{Data-driven solutions for the flux vector field for the natural formulation and all of its variants with a fixed number of sampled points $ N_{g} \times N_{\theta} = 6 \times 6 $ and three different number of finite element meshes. }
	\def\svgwidth{\columnwidth}
	\input{Figures/sol_DNLSP1_plt_over_line_s_fem_mesh.pdf_tex}
	\label{fig:sol_DNLSP1_plt_over_line_s_fem_mesh}
\end{figure}



%\begin{table}[h!]
%	\centering
%	\caption{$L^{2}$-error for the scalar pressure field $\uD_h$, the flux vector field $\sD_h$, and the corresponding dual fields $\lD_h$ and $\mD_h$ for the natural formulation and all its variants, considering a fixed number of sampled points $N_{g} \times N_{\theta} = 6 \times 6$ and three different finite element meshes.}
%	\resizebox{\textwidth}{!}{%
%		\begin{tabular}{lccccc}
%			\toprule
%			\textbf{Formulation} & \textbf{FEM Mesh} & $ \| \uD - \uD_h \|_{L^{2}} $ & $ \| \sD - \sD_h \|_{L^{2}} $ & $ \| \lD - \lD_h \|_{L^{2}} $ & $ \| \mD - \mD_h \|_{L^{2}} $ \\
%			\midrule
%			\multirow{3}{*}{DNLSP1-A} 
%			& 10.0  & 0.029667 & 0.623456 & 0.049219 & 1.264457 \\
%			& 50.0  & 0.028286 & 0.536606 & 0.056152 & 1.073735 \\
%			& 100.0 & 0.028333 & 0.528759 & 0.056566 & 1.057780 \\
%			\midrule
%			\multirow{3}{*}{DNLSP1-Aeqo} 
%			& 10.0  & 0.031100 & 0.469149 & 0.057765 & 0.917430 \\
%			& 50.0  & 0.028398 & 0.506107 & 0.056653 & 1.011805 \\
%			& 100.0 & 0.028369 & 0.512206 & 0.056703 & 1.024314 \\
%			\midrule
%			\multirow{3}{*}{DNLSP1-B} 
%			& 10.0  & 0.059674 & 0.544448 & 0.058580 & 1.011495 \\
%			& 50.0  & 0.030201 & 0.518981 & 0.056638 & 1.034829 \\
%			& 100.0 & 0.028834 & 0.518285 & 0.056706 & 1.035786 \\
%			\midrule
%			\multirow{3}{*}{DNLSP1-C} 
%			& 10.0  & 0.059674 & 0.623456 & 0.049219 & 1.011495 \\
%			& 50.0  & 0.030201 & 0.536606 & 0.056152 & 1.034829 \\
%			& 100.0 & 0.028834 & 0.528759 & 0.056566 & 1.035786 \\
%			\midrule
%			\multirow{3}{*}{DNLSP1-D} 
%			& 10.0  & 0.029667 & 0.544448 & 0.058580 & 1.264457 \\
%			& 50.0  & 0.028286 & 0.518981 & 0.056638 & 1.073735 \\
%			& 100.0 & 0.028333 & 0.518285 & 0.056706 & 1.057780 \\
%			\bottomrule
%		\end{tabular}%
%	}
%	\label{tab:DNLSP1_fem_mesh_errors}
%\end{table}

The same study was conducted for the flux vector field, as summarized in Figure \ref{fig:sol_DNLSP1_plt_over_line_s_fem_mesh}. Again, we observe that the data-driven solutions based on formulations A and C exhibited similar profiles due to the use of a discontinuous Galerkin finite element space. Conversely, the data-driven solutions based on formulations B and D were also similar, as the same finite element space was chosen for that unknown. The data-driven solution for the flux based on the equal-order formulation, however, differs from the others.

Regarding the accuracy of the numerical solution for $ \sD_h $, we conclude that, as with $ \uD_h  $, the convergence is asymptotic up to a certain threshold for formulations A, B, C, and D, with only marginal changes in $ \|\sD - \sD_h \|_{L^{2}} $ when refining the mesh from $ 50 \times 50 $ to $ 100 \times 100 $ elements per coordinate direction. Finally, for the sake of brevity, we do not present the two-dimensional plots and the associated plot over lines for the dual field in this study. 



%--------------------------------------------------------------------------    
%\begin{proposition}
%	For a given finite element mesh with a fixed spatial step $ h $, the data-driven solution $ (\tilde{\sD}_h, \tilde{\eD}_h) $ converges to the reference solution $ (\sD_h^{\text{ref}}, \eD_h^{\text{ref}}) $ in a consistent norm \( \|\cdot\| \) as the number of sampled data points $ N_{m} \rightarrow \infty$ .
%\end{proposition}

%\begin{proof}
%	
%	Let $\sD_h^{\text{ref}}$ and $\eD_h^{\text{ref}}$ be the reference functions, and let $\hat{\sD}_h$ and $\hat{\eD}_h$ be their projections onto the space $\text{DG}_0$. These projections minimize the $L^2$-norm projection errors:
%	
%	\[
%	\| \sD_h^{\text{ref}} - \hat{\sD}_h \|_{L^2(\Omega)} \quad \text{and} \quad \| \eD_h^{\text{ref}} - \hat{\eD}_h \|_{L^2(\Omega)},
%	\]
%	
%	implying that $\hat{\sD}_h$ and $\hat{\eD}_h$ are the closest elementwise constant functions in $\text{DG}_0$ to $\sD_h^{\text{ref}}$ and $\eD_h^{\text{ref}}$, respectively. As $N_{m} \to \infty$, the discrete data set $\mathfrak{D}_m$ becomes dense in the solution space. By definition, a set $\mathfrak{D}_m$ is dense in a solution space if every point in that space can be approximated arbitrarily closely by point from $\mathfrak{D}_m$. Specifically, for any pair $(\hat{\sD}_h, \hat{\eD}_h)$ in the solution space and any $\epsilon > 0$, there exist discrete data pair $\left\{\left(\mathfrak{g}_n, \mathfrak{s}_n\right)\right\}_{n \in \mathcal{N}_m}$ such that:
%	
%	\[
%	\| (\tilde{\sD}_h|_{K_n}, \tilde{\eD}_h|_{K_n}) - (\hat{\sD}_h|_{K_n}, \hat{\eD}_h|_{K_n}) \| < \epsilon.
%	\]
%	
%	Thus, as $N_{m}$ increases, the discrete data $\mathfrak{D}_m$ densely samples the solution space and the optimization problem ensures:
%	
%	\[
%	\| (\tilde{\sD}_h|_{K_n}, \tilde{\eD}_h|_{K_n}) - (\hat{\sD}_h|_{K_n}, \hat{\eD}_h|_{K_n}) \| \to 0 \text{ as } N_{m} \to \infty.
%	\]
%	
%	Thus, we proceed as follows: 
%	
%	\begin{equation*}
%		\begin{aligned}
%			\| (\tilde{\sD}_h, \tilde{\eD}_h) - (\sD_h^{\text{ref}}, \eD_h^{\text{ref}}) \| &=  \| (\tilde{\sD}_h, \tilde{\eD}_h) - (\hat{\sD}_h, \hat{\eD}_h)  +  (\hat{\sD}_h, \hat{\eD}_h) - (\sD_h^{\text{ref}}, \eD_h^{\text{ref}}) \| \\
%			&\leq \| (\tilde{\sD}_h, \tilde{\eD}_h) - (\hat{\sD}_h, \hat{\eD}_h) \| + \| (\hat{\sD}_h, \hat{\eD}_h) - (\sD_h^{\text{ref}}, \eD_h^{\text{ref}}) \|.
%		\end{aligned}
%	\end{equation*}
%	
%	
%	As $N_{m} \to \infty$, the term $\| (\tilde{\sD}_h, \tilde{\eD}_h) - (\hat{\sD}_h, \hat{\eD}_h) \| \to 0$. Since $\| (\hat{\sD}_h, \hat{\eD}_h) - (\sD_h^{\text{ref}}, \eD_h^{\text{ref}}) \|$ is minimized, the difference $\| (\tilde{\sD}_h, \tilde{\eD}_h) - (\sD_h^{\text{ref}}, \eD_h^{\text{ref}}) \|$ also approaches zero.
%	
%\end{proof}
%----------------------------------------------------------------------------

\clearpage
\subsubsection{Convergence of the data-driven formulation}\label{sec:convergence_ddcm}


\begin{figure}[h!]
	\centering
	\caption{Figures (a) to (e) depict the $ L^{2} $-norm convergence error of the scalar pressure field in Darcy's problem, $ \|\uD-\uD_h\|_{L^{2}} $, for each data-driven formulation. Figure (f) presents the reference solution for the scalar pressure field in the manufactured Darcy problem.}
	\def\svgwidth{1,0\linewidth}  
	\input{Figures/DNLSP1_conv_u_L2.pdf_tex}  
	\label{fig:DNLSP1_conv_u_L2}  
\end{figure}

\begin{figure}[h!]
	\centering
	\caption{Figures (a) to (c) depict the $ H^{1} $-norm convergence error of the scalar pressure field in Darcy's problem, $ \|\uD-\uD_h\|_{H^{1}} $, for each data-driven formulation, whose finite element space for $ \uD_h $ consists of $ \text{CG}_{k+1}^{D} $.}
	\def\svgwidth{\columnwidth}  
	\input{Figures/DNLSP1_conv_u_H1.pdf_tex}  
	\label{fig:DNLSP1_conv_u_H1}  
\end{figure}

\begin{figure}[h!]
	\centering
	\caption{Figures (a) to (e) depict the $ L^{2}$-norm convergence error of the flux vector field in Darcy's problem, $ \|\sD-\sD_h\|_{L^{2}} $, for each data-driven formulation. Figure (f) presents the reference solution for the flux vector field in the manufactured Darcy problem.}
	\def\svgwidth{\columnwidth}  
	\input{Figures/DNLSP1_conv_s_L2.pdf_tex}  
	\label{fig:DNLSP1_conv_s_L2}  
\end{figure}

\begin{figure}[h!]
	\centering
	\caption{Figures (a) and (b) depict the $Hdiv$-norm convergence error of the scalar pressure field in Darcy's problem, $ \|\sD-\sD_h\|_{Hdiv} $, for each data-driven formulation, whose finite element space for $ \sD_h $ consists of $ \text{RT}_{k+1} $.}
	\def\svgwidth{\columnwidth}  
	\input{Figures/DNLSP1_conv_s_Hdiv.pdf_tex}  
	\label{fig:DNLSP1_conv_s_Hdiv}  
\end{figure}


\begin{figure}[h!]
	\centering
	\caption{Figures (a) to (e) depict the $ L^{2}$-norm convergence error of scalar dual field in Darcy's problem, $ \|\lambda-\lambda_h\|_{L^{2}} $, for each data-driven formulation.}
	\def\svgwidth{\columnwidth}  
	\input{Figures/DNLSP1_conv_lambda_L2.pdf_tex}  
	\label{fig:DNLSP1_conv_lambda_L2}  
\end{figure}

\begin{figure}[h!]
	\centering
	\caption{Figures (a) to (e) depict the $ L^{2}$-norm convergence error of vector dual field in Darcy's problem, $ \|\boldsymbol{\mu}-\boldsymbol{\mu}_h\|_{L^{2}} $, for each data-driven formulation.}
	\def\svgwidth{\columnwidth}  
	\input{Figures/DNLSP1_conv_mu_L2.pdf_tex}  
	\label{fig:DNLSP1_conv_mu_L2}  
\end{figure}




\clearpage
\subsubsection{Effect of noise on the input data}

So far, we have studied data-driven solutions to the Darcy problem, examining the effects of variables such as the number of finite elements in the mesh, $ (N_{x} \times N_{y}) $, the number of sampled points, $ (N_{m} = N_{g} \times N_{\theta}) $, and the convergence of unknown data fields: the pressure field $ \uD $, the flux vector field $ \sD $, the pressure gradient field $ \eD $, and the dual fields $ \mD $ and $ \lD $ for the natural data-driven formulation and its variants. However, these studies assume a given dataset, $\mathfrak{D}_{m} = \{(\mathfrak{g}_{n}, \mathfrak{s}_{n})\}_{n \in \mathcal{N}_{m}}$, that is noise-free. The objective of this section is to investigate the characteristics of a data-driven solution when a Gaussian noise is introduced into the available data. To this end, we define the standard normal distribution by its mean and standard deviation, which will be denoted by $ \mu $ and $ \sigma $, respectively. Moreover, since the maximum and minimum values of both $(\eD_{h}^{\text{ref}}, \sD_{h}^{\text{ref}})$ are known \textit{a priori} from the manufactured solution, the random noise generated for these data fields undergoes a min-max scaling process to ensure that the noise range aligns with that of the manufactured solution. Mathematically, the noise-added data fields can be expressed as:

\begin{equation}
	\begin{aligned}
		\mathfrak{g}_{\text{noise}} &= \left[\frac{\mathfrak{g}_{\text{gn}} - \min(\mathfrak{g}_{\text{gn}})}{\max(\mathfrak{g}_{\text{gn}}) - \min(\mathfrak{g}_{\text{gn}})} \right] \left[\max(\mathfrak{g}) - \min(\mathfrak{g})\right] + \min(\mathfrak{g}), \\
		\mathfrak{s}_{\text{noise}} &= \left[\frac{\mathfrak{s}_{\text{gn}} - \min(\mathfrak{s}_{\text{gn}})}{\max(\mathfrak{s}_{\text{gn}}) - \min(\mathfrak{s}_{\text{gn}})} \right] \left[\max(\mathfrak{s}) - \min(\mathfrak{s})\right] + \min(\mathfrak{s}),
	\end{aligned}
\end{equation}

where $ \mathfrak{g}_{\text{gn}} $ and $ \mathfrak{s}_{\text{gn}} $ are vectors with components $ g_{i} $ and $ s_{i} $ that are independently drawn from a univariate normal distribution, i.e., $ g_{i}, s_{i} \sim \mathcal{N}(\mu, \sigma) \quad \forall i \in \{1, \dots, d\} $. Finally, to control the magnitude of the introduced noise, we introduce a scalar parameter $ \alpha \in \mathbb{R} $, which scales $ \mathfrak{g}_{\text{noise}} $ and $ \mathfrak{s}_{\text{noise}} $. This parameter, $ \alpha $, will be referred to as the magnitude factor.

To carry out this study, we selected the DNLSP1-A formulation and assumed a Gaussian distribution with zero mean and unit standard deviation. The number of sampled points was set as follows: \( N_{g} = 15 \) with three distinct values for \( N_{\theta} = \{5, 10, 20\} \). Additionally, four values for the magnitude factor were considered: \( \alpha = \{0.05, 0.1, 0.2, 0.4\} \). This range begins with a very low level of noise and progressively increases to a higher level, where certain points may be regarded as outliers in the context of an experimental test. Finally, the finite element mesh was maintained as one hundred elements per coordinate direction, i.e. $ (N_{x}\times N_{y} = 100 \times 100) $.

Figure \ref{fig:compound_noise_u} then presents the data-driven solutions obtained for the scalar pressure field. As shown, the data-driven solution closely approximates the expected analytical solution, even when subjected to a high-intensity noise dataset. However, it is noteworthy that for \( N_{\theta} = 5 \), there is a noticeable deviation between the data-driven pressure profiles and the exact solution for all values of \( \alpha \). This deviation can be attributed to the limited number of sampled points, as discussed in previous sections. Figure \ref{fig:compound_noise_s}, on the other hand, displays the data-driven solutions for the flux vector field. At first glance, Figure \ref{fig:compound_noise_s}(a) might suggest that the Gaussian noise significantly impacts the accuracy of the data-driven solution for the flux vector field. However, a comparison with Figure \ref{fig:DNLSP1_A_plt_over_line_s}(e) reveals that the observed discrepancies between the data-driven solutions and the expected analytical solution are primarily due to the low number of sampled points. So, although the introduction of Gaussian noise seems to affect the flux vector field more than the scalar pressure field, the data-driven solutions still exhibit satisfactory accuracy, even when subjected to a high-intensity noise dataset.


\begin{figure}[h!]
	\centering
	\caption{Figures (a), (f), and (k) display the data-driven solution for the scalar pressure field along the diagonal coordinate based on DNLSP1-A formulation for varying values of $\alpha$ and for each specified value of $N_{\theta}$, in the given order. Figures (b) to (e), (g) to (j), and (l) to (o) present the 2-norm of the flux-gradient relationship for each magnitude factor, corresponding to all considered values of $N_{\theta}$, respectively.}
%	\fontsize{6pt}{10pt}\selectfont 
	\def\svgwidth{\columnwidth}  
	\input{Figures/compound_noise_u.pdf_tex}  
	\label{fig:compound_noise_u}  
\end{figure}

\begin{figure}[h!]
	\centering
	\caption{Figures (a), (f), and (k) display the data-driven solution for the flux vector field along the diagonal coordinate based on DNLSP1-A formulation for varying values of $\alpha$ and for each specified value of $N_{\theta}$, in the given order. Figures (b) to (e), (g) to (j), and (l) to (o) present the two-dimensional flux vector field for each magnitude factor, corresponding to all considered values of $ N_{\theta} $, respectively.  }
	\def\svgwidth{\columnwidth}  
	\input{Figures/compound_noise_s.pdf_tex}  
	\label{fig:compound_noise_s}  
\end{figure}


%---------------------------------------------------------------------------
% I chose to omit the dual fields for the study of the effect of noise, since it does not convey much information.


%Finally, regarding the data-driven solutions for the dual field, i.e. $ \lD $ and $ \mD $, Figures \ref{fig:compound_noise_l} and \ref{fig:compound_noise_m} show the obtained results for varying values of $ \alpha $ and each specified value of $ N_{\theta} $. 

%\begin{figure}[h!]
%	\centering
%	\caption{Data-driven solution for the scalar dual field, $\lD$, based on the DNLSP1-A formulation for varying values of $\alpha$ and each specified value of $N_{\theta}$.}
%	\def\svgwidth{\columnwidth}  
%	\input{Figures/compound_noise_l.pdf_tex}  
%	\label{fig:compound_noise_l}  
%\end{figure}
%
%\begin{figure}[h!]
%	\centering
%	\caption{Data-driven solution for the vector dual field, $\mD$, based on the DNLSP1-A formulation for varying values of $\alpha$ and each specified value of $N_{\theta}$.
%	}
%	\def\svgwidth{\columnwidth}  
%	\input{Figures/compound_noise_m.pdf_tex}  
%	\label{fig:compound_noise_m}  
%\end{figure}
%-----------------------------------------------------------------------------



\clearpage
\subsubsection{A non-linear example}



\clearpage
\subsection{Results for elasticity problem}

We now address the classical elasticity problem. To introduce some novelties in this study, we consider nontrivial geometries both for the domain over which synthetic data are generated and for the one in which the problem is to be solved by the data-driven formulation. For the synthetic data generation process, we again propose numerical experiments that aim to imitate a typical stress-strain experiment performed in the laboratory. In this regard, by means of the GMSH library, we construct a geometric domain that resembles a material specimen and whose dimensions are illustrated in figure \ref{fig:compound_meshes_elastostatic}(a); similarly, for the data-driven domain, we chose to use the classical Cook membrane, figure \ref{fig:compound_meshes_elastostatic}(b), since it is a classical benchmark problem in computational elasticity.

\begin{figure}[h!]
	\centering
	\caption{Figure (a) on the left shows the geometric domain that mimics a material specimen and is used to generate synthetic data. Figure (b) on the right shows the Cook membrane domain where the data-driven problem is to be solved.}
	\def\svgwidth{\columnwidth}
	\input{Figures/compound_meshes_elastostatic.pdf_tex}
	\label{fig:compound_meshes_elastostatic}
\end{figure}


To mathematically formalize both problems, i.e., the synthetic data generation problem and the Cook membrane data-driven problem, we start by denoting the specimen domain as a closed, bounded, and connected subset of the two-dimensional Euclidean space, $ \Omega_{\text{sp}} \subset \mathbb{R}^{2} $, whose boundary can be conveniently defined as: $ \partial \Omega_{\text{sp}} = \left\{\Gamma_{\text{bottom}}^{\text{sp}} \cup \Gamma_{\text{right}}^{\text{sp}} \cup \Gamma_{\text{top}}^{\text{sp}} \cup \Gamma_{\text{left}}^{\text{sp}} \right\} $, where 

\begin{equation*}
	\begin{aligned}
		&\Gamma_{\text{bottom}}^{\text{sp}} && = ~ \left\{ \mathbf{x} \in \Omega_{\text{sp}}\mid x_1 \in \left[ 0, \frac{H}{5}\right], \, x_2 = 0 \right\}\\
		&\Gamma_{\text{top}}^{\text{sp}} && = ~\left\{ \mathbf{x} \in \Omega_{\text{sp}} \mid x_1 \in \left[ 0, \frac{H}{5}\right], \, x_2 = H \right\}\\
		&\Gamma_{\text{left}}^{\text{sp}} && = ~ \left\{ \mathbf{x} \in \Omega_{\text{sp}} \mid x_1 = 0, \, x_2 \in \left[0,H\right] \right\}\\
		&\Gamma_{\text{right}}^{\text{sp}} &&= ~ \bigcup_{i=1}^{5}\Gamma_{i}, \\
	\end{aligned}
\end{equation*}

with each $ \Gamma_{i} $ given by:

{\small
\begin{equation*}
		\begin{aligned}
			&\Gamma_{1} = \left\{\mathbf{x} \in \Omega_{\text{sp}} \mid x_{1} = \frac{H}{5}, \ x_{2} \in \left[0,\frac{H}{5}\right] \right\} \\
			&\Gamma_{2} = \left\{\mathbf{x} \in \Omega_{\text{sp}} \mid x_{1} = \frac{H}{5}, \ x_{2} \in \left[\frac{4H}{5}, H\right] \right\} \\
			&\Gamma_{3} = \left\{\mathbf{x} \in \Omega_{\text{sp}} \mid x_{1} = \frac{H}{10}, \ x_{2} \in \left[\frac{3H}{10}, \frac{7H}{10}\right] \right\} \\
			&\Gamma_{4} = \left\{\mathbf{x} \in \Omega_{\text{sp}} \mid x_{1} \leq \frac{H}{5}, \ x_{2} \leq \frac{3H}{10}, \ \left(x_{1}-\frac{H}{5}\right)^2 + \left(x_{2}-\frac{3H}{10}\right)^2 = \left(\frac{H}{10}\right)^2\right\} \\
			&\Gamma_{5} = \left\{\mathbf{x} \in\Omega_{\text{sp}} \mid x_{1} \leq \frac{H}{5}, \ x_{2} \geq \frac{4H}{5}, \ \left(x_{1}-\frac{H}{5}\right)^2 + \left(x_{2}-\frac{H}{5}\right)^2 = \left(\frac{H}{10}\right)^2\right\}
		\end{aligned}
\end{equation*}
}

Similarly, we also consider the Cook membrane domain a closed, bounded, and connected subset of the two-dimensional Euclidean space, $ \Omega_{\text{ck}} \subset \mathbb{R}^{2} $, whose boundary can be defined as $ \partial \Omega_{\text{ck}} = \left\{\Gamma_{\text{bottom}}^{\text{ck}} \cup \Gamma_{\text{right}}^{\text{ck}} \cup \Gamma_{\text{top}}^{\text{ck}} \cup \Gamma_{\text{left}}^{\text{ck}} \right\} $, where

\begin{equation*}
	\begin{aligned}
		&\Gamma_{\text{bottom}}^{\text{ck}} && = ~ \left\{ \mathbf{x} \in \Omega_{\text{ck}} \mid x_{2} = \frac{11}{12}x_{1} \right\}\\
		&\Gamma_{\text{top}}^{\text{ck}} && = ~\left\{ \mathbf{x} \in \Omega_{\text{ck}} \mid x_{2} = \frac{1}{3}x_{1}+ \frac{11B}{12}\right\}\\
		&\Gamma_{\text{left}}^{\text{ck}} && = ~ \left\{ \mathbf{x} \in \Omega_{\text{ck}} \mid x_1 = 0, \, x_2 \in \left[0,H\right] \right\}\\
		&\Gamma_{\text{right}}^{\text{ck}} &&= ~ \left\{ \mathbf{x}  \in \Omega_{\text{ck}} \mid x_{1} = B, \ x_{2} \in \left[\frac{11B}{12},\frac{5B}{4}\right]\right\} \\
	\end{aligned}
\end{equation*}

Once we have formally defined each domain, we must specify the process by which we generate the synthetic pairs $\mathfrak{D}_{m} = \left\{\mathfrak{e}_{n},\mathfrak{s}_{n}\right\}_{n \in \mathcal{N}_{m}}$, which are now second-order tensors, denoting the symmetric gradient of the displacement field and the stress field, respectively, and which will ultimately be used to compute the functions $\tilde{\eE}_h$ and $\tilde{\sE}_h$, as we have already described in section \ref{sec:datadist}. To do this, we consider two independent numerical experiments: in the first, we subject the specimen domain to a \textbf{uniaxial} stress test, as illustrated in figure \ref{fig:synthetic_data_elastostatic}(a); in the second, we subject the specimen domain to a \textbf{shear} stress test, as shown in figure \ref{fig:synthetic_data_elastostatic}(b). The boundary conditions to be imposed to perform both the uniaxial and shear stress tests in $ \Omega_{\text{sp}} $ are defined by equation \ref{eq:bcs_elasticity_synth_data}:

{\small
\begin{equation}\label{eq:bcs_elasticity_synth_data}
	\textbf{Uniaxial:} \begin{cases}
		\begin{aligned}
			&u_{y} =-t,~&&\text{on}~\Gamma_{\text{top}}^{\text{sp}}	\\
			&u_{y} = 0,~&&\text{on}~\Gamma_{\text{bottom}}^{\text{sp}} \\
			&\sE \cdot \mathbf{\breve{n}} = [0,0]^{\top},~&&\text{on}~\Gamma_{\text{right}}^{\text{sp}} \\
			&u_{x} = 0,~&&\text{on}~\Gamma_{\text{left}}^{\text{sp}} \\
		\end{aligned}
	\end{cases} 
	\textbf{Shear:} \begin{cases}
		\begin{aligned}
			&u_{x} =-t,~&&\text{on}~\Gamma_{\text{top}}^{\text{sp}}	\\
			&\mathbf{u} = [0,0]^{\top},~&&\text{on}~\Gamma_{\text{bottom}}^{\text{sp}} \\
			&\sE \cdot \mathbf{\breve{n}} = [0,0]^{\top},~&&\text{on}~\Gamma_{\text{right}}^{\text{sp}}\\
			&\sE \cdot \mathbf{\breve{n}} = [0,0]^{\top},~&&\text{on}~\Gamma_{\text{left}}^{\text{sp}}
		\end{aligned}
	\end{cases} 
\end{equation}
}

where $ t $ is the parametric value used to impose different values of vertical displacement in the uniaxial stress test and horizontal displacement in the shear stress test. 


\begin{figure}[h!]
	\centering
	\caption{Figure (a) at the top shows numerical solution of the displacement field for the uniaxial stress test on the specimen geometry considering different values of vertical displacement on $\Gamma_{\text{top}}^{\text{sp}}$ and parametrized by $t$. Figure (b) at the bottom presents the numerical solution fo the displacement field for the shear stress test on the same geometry, with different values of horizontal displacement on $\Gamma_{\text{top}}^{\text{sp}}$ parametrized by $t$. In both figures, the points where the stress-strain measurements are evaluated are highlighted, simulating an extensometer setup. }
	\def\svgwidth{\columnwidth}  
	\input{Figures/synthetic_data_elastostatic.pdf_tex}  
	\label{fig:synthetic_data_elastostatic}  
\end{figure}

We now move on to the Cook membrane data-driven problem. For this study, we consider an embedded Cook membrane on $\Gamma_{\text{left}}^{\text{ck}}$ and subjected to an vertical shear force on $\Gamma_{\text{right}}^{\text{ck}}$. Mathematically, the imposed boundary conditions are given by equation \ref{eq:bcs_cook_membrane}.

\begin{equation}\label{eq:bcs_cook_membrane}
	\begin{cases}
		\begin{aligned}
		&\mathbf{u} = [0,0]^{\top}&&\text{on}~\Gamma_{\text{top}}^{\text{sp}} \\
		&\sE \cdot \mathbf{\breve{n}} = [0,0]^{\top},~&&\text{on}~\Gamma_{\text{bottom}}^{\text{sp}}\\
		&\sE \cdot \mathbf{\breve{n}} = \left[0,\mathcal{F}_{y}\right]^{\top},~&&\text{on}~\Gamma_{\text{right}}^{\text{sp}} \\
		&\sE \cdot \mathbf{\breve{n}} = [0,0]^{\top},~&&\text{on}~\Gamma_{\text{left}}^{\text{sp}}\\
		\end{aligned}
	\end{cases}
\end{equation}

To perform the numerical experiment, we set the dimensions of the specimen and the Cook membrane to be $ H=1 $ and $ B=1 $, respectively (Figure \ref{fig:compound_meshes_elastostatic}). Regarding the finite element meshes, the characteristic length of the specimen was set to $ l_c = 0.05 $ meanwhile a mesh comprising $ 20.000 $ triangular elements was adopted to discretize the Cook membrane geometry. Concerning the synthetic data generation process, we chose to sample a set of $ 10 $ points across the specimen geometry (Figure \ref{fig:synthetic_data_elastostatic}), each of them obtained by imposing $ N_{g}=100 $ equally spaced displacement values on $ \Gamma_{\text{top}}^{\text{sp}} $, whose range encompasses the minimum and maximum displacements of the target solution. 

\begin{figure}[h!]
	\centering
	\caption{Figures (a) and (b) compare the numerical displacement field in the Cook membrane obtained using the classical primal finite element method and the data-driven approach. Likewise, Figures (c) and (d) present a comparison of the numerical stress field in the Cook membrane, highlighting the differences between the classical finite element formulation and the data-based formulation. }
	\def\svgwidth{\columnwidth}  
	\input{Figures/ddcm_cook_membrane.pdf_tex}  
	\label{fig:ddcm_cook_membrane}  
\end{figure}

Additionally, the obtained pairs $ \left\{\mathfrak{e},\mathfrak{s}\right\} $ were rotated by $ N_{\theta} = 100 $ angles uniformly distributed in the interval $ [0, 2\pi) $. Both processes were conducted for each of the tests, i.e., uniaxial tension and shear. The collected results are, then, summarize in Figure \ref{fig:ddcm_cook_membrane}. As can be observed, the, the data-driven solution greatly approximates to the expected reference solution for both displacement and stress fields. The experiment evidence the robustness of the data-driven approach, in the sense that it can handle with different problems even though they do not share the same geometry characteristics


\subsubsection{Convergence to a manufactured solution}

(Preliminary results: sampling from the reference problem)

To establish the convergence study for the elasticity problem following the same model we adopted to Darcy's problem, an manufactured solution is proposed on a regular Cartesian two-dimensional domain, $ \Omega \subset \mathbb{R}^{2} \mid \Omega = [0,1] \times [0,1] $, with pure Dirichlet boundary conditions. To this end, we consider the following manufactured displacement field

\begin{equation*}
	\begin{aligned}
		\mathbf{u}\colon \Omega \rightarrow \mathbb{R}^{2}, \quad \mathbf{u} = \left[\cos(\pi x)\cos(\pi y),\cos(\pi x)\cos(\pi y) \right]^{\top}	
	\end{aligned}
\end{equation*}

from which one can derive the associated stress, $ \sE $, symmetric gradient of displacement field, $ \eE $, and the body force, $ \mathbf{f} $

\begin{equation*}
	\begin{aligned}
		&\sE \colon \Omega \rightarrow \mathbb{R}^{2\times2}, \quad &&\sE = \begin{bmatrix}
			-\pi\left[\alpha\sin(\pi x) + \beta\sin(\pi y)\right] & 0 \\
			0 &-\pi\left[\alpha\sin(\pi y) + \beta\sin(\pi x)\right]
		\end{bmatrix}
		 \\
	 	&\eE \colon \Omega \rightarrow \mathbb{R}^{2\times2}, \quad &&\eE = \begin{bmatrix}
	 		-\pi \sin(\pi x) & 0 \\
	 		0 &-\pi \sin(\pi y)
	 	\end{bmatrix}
 		\\
 		&\mathbf{f} \colon \Omega \rightarrow \mathbb{R}^{2}, \quad &&\mathbf{f} = \left[\pi^2\alpha\cos(\pi x), \pi^2\alpha\cos(\pi y) \right]^{\top}
	\end{aligned}
\end{equation*}


being $ \alpha = (\lambda_L + 2\mu_L) $ and $ \beta = \lambda_L $. 




\begin{figure}[h!]
	\centering
	\caption{... }
	\def\svgwidth{\columnwidth}  
	\input{Figures/ddcm_convergence_elasticity.pdf_tex}  
	\label{fig:ddcm_convergence_elasticity}  
\end{figure}


\cleardoublepage

